\subsection{Re-folding devices and collars}

\begin{proof}[Proof of Lemma~\ref{lem:refold}]
(a) After the class-uniform permutation the branch is uniform on its class: on both points $c\pm D$ of its pair and on all $m_{t-1}$ slots. The round moves the point by the
label's translation and the device shifts it by a fixed integer, so the new distance is $X=|D+\sigma\delta_t|$ with $\sigma$ uniform, and the slot is still uniform. The
source pair at distance $X$ has rank $j=X-\frac12$; its $2m_{t-1}$ units have ranks $2jm_{t-1}+2k+s$ (slot $k$, point $s\in\{0,1\}$), and target pair $i$ holds the ranks
$[2m_ti,2m_t(i+1))$. A unit of slot $k$ goes below target pair $j$ if and only if $2k+s<2j(m_t-m_{t-1})$, that is, if and only if $k<j(m_t-m_{t-1})$, for both points. So
if $j\gamma_t\le1$ a fraction exactly $j\gamma_t$ of the slots goes to pair $j-1$ and the rest to pair $j$, since pair $j$ ends at $2m_t(j+1)\ge2jm_{t-1}+2m_{t-1}$; and if
$j\gamma_t>1$ all units go below pair $j$.

(b) Put $\lambda=\sqrt{\gamma L/2}/\delta_{\max}\le1/(2\delta_{\max}^2)\le\frac12$ and $D_1=\delta_{\max}+\frac12+2\lambda\delta_{\max}^2/\gamma$. For $D\ge D_1$, $X=D+\sigma\delta$ and
$q:=\min(1,(X-\frac12)\gamma_t)\ge\min(1,(D-\delta_{\max}-\frac12)\gamma)\ge2\lambda\delta_{\max}^2$, and in both cases of (a) $\E[e^{\lambda D'}\mid X]\le e^{\lambda X}(1-q(1-e^{-\lambda}))$. With
$1-e^{-\lambda}\ge\lambda/2$ and $\cosh(\lambda\delta)\le e^{\lambda^2\delta^2/2}$, $\E[e^{\lambda D'}\mid D]\le e^{\lambda D}\exp(\lambda^2\delta_{\max}^2/2-q\lambda/2)\le e^{\lambda D}e^{-\lambda^2\delta_{\max}^2/2}
\le(1-\theta)e^{\lambda D}$ with $\theta=\lambda^2\delta_{\max}^2/4$, as $\lambda^2\delta_{\max}^2/2\le1$. For $D<D_1$, $D'\le X<D_1+\delta_{\max}$, so $e^{\lambda D'}\le K:=e^{\lambda(D_1+\delta_{\max})}$.
Hence $V_t=\E e^{\lambda D_t}$ satisfies $V_{t+1}\le(1-\theta)V_t+K$ and $V_t\le e^{\lambda R_0}+K/\theta$. By Markov's inequality and a union over $t\le n$,
$\Pr(\exists t:D_t\ge y)\le n(e^{\lambda R_0}+K/\theta)e^{-\lambda y}$. At $y=R_0+2\delta_{\max}+\frac12+2L/\lambda$: $L/\lambda=\delta_{\max}\sqrt{2L/\gamma}=2\lambda\delta_{\max}^2/\gamma$, so
$y-D_1-\delta_{\max}=R_0+L/\lambda$, and the bound is at most $ne^{-2L}+n(8/(\gamma L))e^{-L}\le ne^{-L}(1+8/(\gamma L))\le\eta$, because $ne^{-L}=\eta/(1+8/\gamma)$ and
$L\ge1$. Finally $2L/\lambda=2\delta_{\max}\sqrt{2L/\gamma}$.
\end{proof}

\begin{proof}[Proof of Theorem~\ref{thm:fibred}]
\emph{The device.} The memory holds the points of $\Omega_W$, $k$ slot registers, one per coordinate, with $m^{(i)}_t$ slots in use at time $t$, $m_0\ge\max(n,1/\gamma)$ and
$m_t=\lceil(1+\gamma)m_{t-1}\rceil$, an ancilla $\C^{m'}$ with a port, and a garbage flag. The windows $\Omega_t$ are the units (a point of $\pi^{-1}(W_x)$ and slots) at time
$t$, and $\Omega_0$ is one pair in each coordinate ($R_0=\frac12$) times the fibre times the slots. The classes are the products of the coordinate pairs $\{c_i\pm D_i\}$, the
fibre and the slots. Round $t$ applies an independent uniformly random permutation of each class, fixed once and for all before any observer; Haar unitaries on the point
blocks $\{\text{unit}\}\otimes\C^{m'}$; the round, the gadget unitary of the regular action on $\Omega_W$ completed to a unitary of the memory; the shift by $-r_0$ and, in each
coordinate, the re-folding injection of Lemma~\ref{lem:refold} by distance, with a fixed bijection between the fibres over the old and the new coordinates; sends units
without a target to garbage; and applies Haar unitaries on the new point blocks and exports. Every operation other than the round acts on the memory alone, and on the window
the round is the gadget unitary of the regular action tensored with the identity on slots and ancillas.

\emph{Exactness and the class chain.} Every point of the window is triangle-good and collision-free, since the action is the regular action of a group in which the
relations hold and the labels are distinct; so the round is a free permutation round on the window. By the fibre hypothesis a uniform point of the fibre over $v$ is mapped
to a uniform point of the fibre over $v-\pi(g)$; the class permutation makes the signs and slots of the coordinates independent and uniform; so the class-uniformized chain
is the product of $k$ independent chains of Lemma~\ref{lem:refold}(a), with steps $\pi_i(g_t)-r_{0,i}$. Lemma~\ref{lem:refold}(b) with $\eta/k$ per coordinate gives
$\max_y\kappa_{\rm class}(y)\le\eta$; a coordinate with $\delta_i=0$ never moves out.

\emph{Discarded weight and error.} Put $z=\min(\epsilon^2/(64n),1/(4n^3))$ and choose $|\Omega_0|$, through $m_0$, with $|\Omega_0|\ge n(n\ln r+\ln2)/(2z^2)$, so that
$\log m_0=O(\log(n/\epsilon))$. By Lemma~\ref{lem:classperm} there are class permutations, fixed from then on, with $\max_y\kappa_\beta(y)\le\eta+z$. Theorem~\ref{thm:pointblock}
with $\zeta=z$, its Haar choices made after the permutations, gives for every observer $\omega_{\rm ideal}\le\eta+z\le\epsilon^2/(32n)$ and $\omega\le\eta+2z\le3/(4n^3)$. So the
error is at most $(1+\sqrt n)\sqrt{\epsilon^2/(32n)+\epsilon^2/(16n)}\le0.53\epsilon$ for $n\ge2$, and the memory mass outside the window at the start of every round is at most
$n\omega\le3/(4n^2)$.

\emph{Memory.} $Q\le\log|\Omega_W|+\sum_i\log m^{(i)}_n+\log m'+O(1)$, with $\log m^{(i)}_n\le\log m_0+n\log(1+\gamma+1/m_0)\le\log m_0+n\log(1+\gamma)+1/\ln2$ since $m_0\ge n$,
$\log m_0=O(\log(n/\epsilon))$, and $\log m'=O_r(\log(n/\epsilon)+\log\log|\Omega_W|)$ by Theorem~\ref{thm:pointblock}.

For $\Z^k$ the fibres are trivial and the slot words move each coordinate by a bounded amount, so $|\Omega_W|\le\prod_i(2x_i+O(1))$; with $\gamma=1/n$ the condition
$\gamma L\le1/(2\delta_{\max}^2)$ holds for $n\ge n_0$, $x_i=O(\delta_i\sqrt{nL})$, and $Q=O_{k,r}(\log(n/\epsilon))$.
\end{proof}

\begin{proof}[Proof of Proposition~\ref{prop:collar}]
(a) $x\in L_j\subseteq F_j$ gives $f(g)x\in F_{j+1}$, and $x\notin F_{j-1}$ with $f(g)^{-1}F_{j-2}\subseteq F_{j-1}$ gives $f(g)x\notin F_{j-2}$. The points of $F_j$ sent outside $F_j$
all lie in $L_j$, since $f(g)F_{j-1}\subseteq F_j$, and there are $|f(g)F_j\setminus F_j|=\Phi_g(j)$ of them. For the down moves, $F_{j-1}\cap f(g)L_j=F_{j-1}\setminus f(g)F_{j-1}$: a
point $y$ of $F_{j-1}$ outside $f(g)F_{j-1}$ has $f(g)^{-1}y\in F_j\setminus F_{j-1}=L_j$, and conversely; and $|F_{j-1}\setminus f(g)F_{j-1}|=|f(g)F_{j-1}\setminus F_{j-1}|=\Phi_g(j-1)$
since $|f(g)F_{j-1}|=|F_{j-1}|$. So $s_jp_+(j)=\Phi_g(j)=s_{j+1}p_-(j+1)$.

(b) \emph{Re-fold.} Ranks are kept and capacities grow ($V_jm_t\ge V_jm_{t-1}$), so no unit moves outward. A unit of shell $X\ge1$ with slot $k$ and index $\mathrm{idx}<s_X$ goes
to a shell at most $X-1$ if and only if $ks_X+\mathrm{idx}<V_{X-1}(m_t-m_{t-1})$; the fraction of slots doing so is at least
$\min(1,V_{X-1}\gamma_t/s_X)-1/m_{t-1}\ge\Gamma':=\min(1,\gamma|F_0|/(2s))-1/m_0$, since $V_{X-1}\ge|F_0|$, $s_X\le2s$, $\gamma_t\ge\gamma$ and $m_{t-1}\ge m_0$. Slots are uniform
and independent of the point, so, given the past, the unit drops by at least one shell with probability at least $\Gamma'$. Reaching the garbage needs shell $K+1$ after a
round, hence shell $K$ before it.

\emph{Drift.} Let $\lambda:=\min(1,\Gamma'/2-\zeta)$. For $D\ge2$ the move $X-D\in\{-1,0,1\}$ has mean $(\Phi_g(D)-\Phi_g(D-1))/s_D\le\zeta$ by (a), so
$\E e^{\lambda(X-D)}\le\exp(\lambda\zeta+\lambda^2/2)$ (Hoeffding's lemma), and $\E[e^{\lambda D'}\mid X]\le e^{\lambda X}(1-\Gamma'(1-e^{-\lambda}))\le e^{\lambda X}\exp(-\Gamma'(\lambda-\lambda^2/2))$.
Since $\lambda\le1$ and $\lambda\le\Gamma'/2-\zeta$, $\Gamma'(1-\lambda/2)\ge\Gamma'/2\ge\zeta+\lambda$, so the exponent $\lambda\zeta+\lambda^2/2-\Gamma'\lambda(1-\lambda/2)$ is at most
$-\lambda^2/2$ and $\E[e^{\lambda D'}\mid\text{past}]\le(1-\lambda^2/4)e^{\lambda D}$. For $D\le1$, $e^{\lambda D'}\le e^{2\lambda}$.

\emph{Kill.} With $U_t:=\E[e^{\lambda D_t}\mathbf 1\{\tau\ge t\}]$, $\tau$ the first visit to shell $K$: $U_t\le(1-\lambda^2/4)U_{t-1}+e^{2\lambda}$, so $U_t\le1+4e^{2\lambda}/\lambda^2\le1+30/\lambda^2$, and
$\Pr(\tau\le n)\le\sum_{t\le n}\Pr(\tau\ge t,D_t\ge K)\le n(1+30/\lambda^2)e^{-\lambda K}$.

\emph{Scales.} $|F_K\setminus F_0|=\sum_{1\le j\le K}s_j\ge Ks$, so $|F_0|\ge Ks$. With $\gamma$ as stated, $\gamma|F_0|/(2s)=2\zeta+4L_c/K+1/m_0\le6L_c/K+\frac14\le1$, so
$\Gamma'=2\zeta+4L_c/K$ and $\lambda=2L_c/K\le\frac14$. The kill is at most $n(1+30K^2/(4L_c^2))e^{-2L_c}\le n(1+8K^2)e^{-2L_c}=\eta e^{-L_c}\le\eta$, since $L_c\ge1$. And
$n\gamma\le(2n/K)(2L_c/K+4L_c/K+1/m_0)\le12nL_c/K^2+2/K$, as $m_0\ge n$, so the slot register costs $\log m_0+n\gamma/\ln2+O(1)\le\log m_0+17.4nL_c/K^2+O(1)$. The rest is the proof of
Theorem~\ref{thm:fibred}, with the core $F_0$ times the initial slots as $\Omega_0$ and $\eta=\min(\epsilon^2/(64n),1/(4n^3))$.
\end{proof}

\begin{proof}[Proof of Theorem~\ref{thm:collar-condition}]
Put $v:=2^{2S+4}$ and take the collar of (a), of depth $K$. If $K\ge8L_c$, with $L_c$ at this depth, Proposition~\ref{prop:collar}(b) applies, since $\zeta\le1/K\le L_c/K$, and gives
$\log\max_t|\Omega_t|\le\kappa(2S+5)(\log(2S+6))^\kappa+17.4nL_c\kappa(\log(2S+6))^\kappa\Lambda(2^{2S+4})+O_r(\log(n/\epsilon)+\log\log|F_{K+1}|)$, and
Theorem~\ref{thm:FK}(ii) gives $n\Lambda(2^{2S+4})\le8C_\mu\,n\,a$. Otherwise let $K_0\ge16$ be the least integer with $K_0\ge8\ln(n(1+8K_0^2)/\eta)$, and take the collar of (b)
for this $K_0$, of depth $K'\ge K_0$. The function $K\mapsto K-8\ln(n(1+8K^2)/\eta)$ increases for $K\ge16$, so $K'\ge8L_c(K')$ and Proposition~\ref{prop:collar}(b) applies: its size
term is at most $\kappa K_0^\kappa=\mathrm{polylog}(n/\epsilon)$, and its $\gamma$-term is $17.4nL_c/K'^2\le17.4nL_c/K^2\le17.4nL_c\kappa(\log(2S+6))^\kappa\Lambda(2^{2S+4})$, since
$K'\ge K_0>K$, which is bounded as before. In both cases $L_c=O(\log(n/\epsilon)+\log K')$ and $\log K'\le\log|F_{K'+1}|$, so all factors are polylogarithmic in $n/\epsilon$ and
$2+S$; (B1) and (B3) hold by the construction of Proposition~\ref{prop:collar}(b). For the regular action of a group of exponential growth, step 2 of the proof of
Proposition~\ref{prop:growth} gives $\Lambda(v)\ge c/(\log v)^2$; at $K\asymp\log v/\log\log v$ a collar of size $e^{O(K\log K)}$ has $\log|F_{K+1}|=O(\log v)$ and
$K^{-2}\le C(\log\log v)^2\Lambda(v)$, which is (a), and (b) is immediate.
\end{proof}

\paragraph{Houghton groups.} For $q\ge3$, $H_q$ acts on $X_q=\{1,\dots,q\}\times\{1,2,\dots\}$ by permutations that are eventually translations on each ray; $t(g)\in\Z^q$, with
$\sum_it_i(g)=0$, is the eventual translation of $g$ on the rays, a homomorphism with kernel the finitary symmetric group $\FSym(X_q)$. The generators $g_j$ ($2\le j\le q$)
have $t(g_j)=e_j-e_1$: $g_j$ moves ray $1$ one step towards the origin, sends $(1,1)$ to $(j,1)$, and moves ray $j$ one step outwards. Put
$s(\vec m)=g_2^{m_2}\circ\cdots\circ g_q^{m_q}$ for $\vec m\in\Z^{q-1}$; every element is uniquely $\sigma\circ s(\vec m)$ with $\sigma\in\FSym(X_q)$, since $x\circ s(\vec m)^{-1}$ has zero
translation for the $\vec m$ read off $t(x)$. The \emph{defect radius} of $g$ is $\varrho(g)=\min\{R\ge1:g(i,k)=(i,k+t_i(g))$ for all $i$ and all $k\ge R\}$.

\begin{lemma}[Cocycles of Houghton groups]\label{lem:houghton-cocycle}
\begin{enumerate}[label=(\alph*),itemsep=1pt]
\item $\varrho(g\circ h)\le\max(\varrho(h),\varrho(g)+\nrm{t(h)}_\infty)$ and $\varrho(g^{-1})\le\varrho(g)+\nrm{t(g)}_\infty$, and $t(c)=0$ implies $c\in\Sym(X_{<\varrho(c)})$, where
$X_{<R}=\{(i,k):k<R\}$ has $q(R-1)$ points.
\item Let $w$ be an element with $\nrm{t(w)}_\infty\le\ell$ and $\varrho(w^{-1})\le\varrho_0$, and let $\pi(w)$ be the $\Z^{q-1}$-part of $t(w)$. For $x=\sigma\circ s(\vec m)$ with
$|m_j-\frac12|\le x_0$ for all $j$, $x\circ w^{-1}=(\sigma\circ c)\circ s(\vec m-\pi(w))$ with $c\in\Sym(X_{<R_\star})$, $R_\star:=\varrho_0+\ell+4+7(q-1)(x_0+1+\ell)$, and $c$ depends
only on $\vec m$ and $w$.
\end{enumerate}
\end{lemma}

\begin{proof}
(a) For $k\ge\max(\varrho(h),\varrho(g)+\nrm{t(h)}_\infty)$, $h(i,k)=(i,k+t_i(h))$ with $k+t_i(h)\ge\varrho(g)$, so $g(h(i,k))=(i,k+t_i(h)+t_i(g))$. For
$k\ge\varrho(g)+\nrm{t(g)}_\infty$ put $k':=k-t_i(g)\ge\varrho(g)$: then $g(i,k')=(i,k)$, so $g^{-1}(i,k)=(i,k-t_i(g))$. If $t(c)=0$, $c$ fixes every $(i,k)$ with $k\ge\varrho(c)$, hence
permutes the finite complement $X_{<\varrho(c)}$.

(b) $x\circ w^{-1}=\sigma\circ s(\vec m)\circ w^{-1}=\sigma\circ c\circ s(\vec m')$ with $\vec m':=\vec m-\pi(w)$ and $c:=s(\vec m)\circ w^{-1}\circ s(\vec m')^{-1}$, which has zero translation and
depends only on $\vec m$ and $w$. Since $\varrho(g_j^{\pm1})=2$ and each factor changes each coordinate of $t$ by at most one, (a) gives $\varrho(s(\vec m))\le2+|\vec m|_1$ and
$\varrho(s(\vec m)^{-1})\le2+2|\vec m|_1$. With $h_1:=s(\vec m')^{-1}$ and $h_2:=w^{-1}\circ h_1$: $\varrho(h_1)\le2+2|\vec m'|_1$, $\nrm{t(h_1)}_\infty\le|\vec m'|_1$,
$\varrho(h_2)\le\max(\varrho(h_1),\varrho_0+|\vec m'|_1)$, $\nrm{t(h_2)}_\infty\le\ell+|\vec m'|_1$, and $\varrho(c)=\varrho(s(\vec m)\circ h_2)\le\max(\varrho(h_2),2+|\vec m|_1+\ell+|\vec m'|_1)$. In
the box, $|\vec m|_1\le(q-1)(x_0+1)$ and $|\vec m'|_1\le(q-1)(x_0+1+\ell)$, so $\varrho(c)\le R_\star$ and $c\in\Sym(X_{<R_\star})$ by (a).
\end{proof}

\begin{proof}[Proof of Theorem~\ref{thm:houghton-groups}]
Use the right regular action $f(g)x=x\circ g^{-1}$, a homomorphism, with every point triangle-good and collision-free. Let $\ell$ bound $\nrm{t(w)}_\infty$ and $\varrho_0$ bound
$\varrho(w^{-1})$ over the slot words. Given the window half-width $x_0$, put $\Omega_B:=\{\sigma\circ s(\vec m):\sigma\in\Sym(X_{<R_\star}),\ \vec m\in B_{\rm out}\}$ with $B_{\rm out}$ the box
of half-width $x_0+\ell$ around $\frac12\cdot\mathbf 1$, and $\pi(\sigma\circ s(\vec m)):=\vec m$. The fibres have the common size $(q(R_\star-1))!$, and by
Lemma~\ref{lem:houghton-cocycle}(b) every slot word maps the fibre over $\vec m$, for $\vec m$ in the box $B$ of half-width $x_0$, bijectively onto the fibre over $\vec m-\pi(w)$,
by $\sigma\mapsto\sigma\circ c$. So Theorem~\ref{thm:fibred} applies with $k=q-1$, $r_0=0$ and $\delta_i\le\ell$, at the window half-width
$x_0=1+2\ell+2\ell\sqrt{2L/\gamma}$, that is $\gamma=8\ell^2L/(x_0-1-2\ell)^2$; its condition $\gamma L\le1/(2\ell^2)$ holds as soon as $x_0\ge1+2\ell+4\ell^2L$. Its memory is
$\log((q(R_\star-1))!)+(q-1)O(\log x_0)+(q-1)n\log(1+\gamma)+O(\log(n/\epsilon))=O(q^2x_0\log(qx_0)+qn\ell^2L/x_0^2)+O(\log(n/\epsilon))$. Take
$x_0=\max\bigl((n\ell^2L/(q\log n))^{1/3},\,1+2\ell+4\ell^2L\bigr)$: the memory is $O(q^{5/3}(\ell^2L)^{1/3}n^{1/3}(\log n)^{2/3}+q^2\ell^2L\log(q\ell L))+O(\log(n/\epsilon))$.

$H_q$ is elementary amenable, being an extension of $\FSym(X_q)$ by $\Z^{q-1}$, and so is the subgroup generated by the labels. By Chou's
theorem~\cite[Theorem~3.2]{Chou1980} that subgroup has polynomial or exponential growth. In the first case it is virtually nilpotent, hence residually finite, and its finite
quotients give a finite-dimensional representation of $P_\Phi$ that separates the labels, so Corollary~\ref{cor:map} applies. In the second case
Proposition~\ref{prop:growth}, on the regular action, gives $S+n\,a\ge c_qn^{1/3}-C_q$ for every bath; since $L=O(\log(n/\epsilon))$, both terms of the memory are at most
$K\log^2(n/\epsilon)(1+S+n\,a)$ for $n\ge n_0$.
\end{proof}

\begin{proof}[Proof of Proposition~\ref{prop:ball}]
(a) With point classes no permutation is needed, and the branch of a label sequence at time $t$ is a point $f(w)x_0$ with $|w|\le t$, inside $\Omega_t$; so the ideal killed
chain is never killed, and Theorem~\ref{thm:pointblock} gives error at most $(1+\sqrt n)\sqrt{\epsilon^2/(16n)}\le0.43\epsilon$ and memory
$\log|\Omega_n|+O_r(\log(n/\epsilon))$. The ball of radius $n$ in the labels and their inverses has at most $(2r+1)^n$ points.

(b) Let $A$ be finite and $r_A:=12(\lfloor\log\log|A|\rfloor+2)$, so that $|A|<2^{2^{\lfloor r_A/12\rfloor-1}}$. By~\cite[Proposition~7.3]{DouglasMghirbiC}, some generator
$s\in\{p,q,t,c,a\}$ of $\Gamma$ has $|As\setminus A|>|A|/r_A$; the generators are labels of the lamp channel, and for the right regular action $f(s)A=As^{-1}$ with
$|As^{-1}\setminus A|=|A\setminus As|=|As\setminus A|$. So the label boundary $|\partial A|=\sum_{s\in S_G}|f(s)A\setminus A|$ exceeds $|A|/\Psi(|A|)$ with
$\Psi(v):=12(\log\log v+2)\ge r_A$. The co-area and Cauchy--Schwarz steps of the proof of Proposition~\ref{prop:growth}, with $|\partial A_t|\ge|A_t|/\Psi(v)$ for the level sets of
$\varphi^2$, give $\mathcal E_{S_G}(\varphi)\ge\nrm\varphi^2/(|S_G|\Psi(v)^2)$ and $\Lambda(v)\ge1/(4r\Psi(v)^2)$. By Theorem~\ref{thm:FK}(ii), $a\ge c/(\log(2S+4))^2$ for every bath on
the regular action, so $S+n\,a\ge c'n/(\log n)^2$: if $S\ge n/(\log n)^2$ this is immediate, and otherwise $\log(2S+4)=O(\log n)$. So the window of (a), of at most $n\log(2r+1)$
bits, is at most $(\log(n/\epsilon))^3(1+S+n\,a)$ for $n\ge n_0$.
\end{proof}

\subsection{Leaky biased rounds}

\begin{proof}[Proof of Theorem~\ref{thm:perm-leaky}]
\begin{enumerate}[label=\arabic*.,itemsep=1pt]
\item By Lemma~\ref{lem:perm-leak}, $\nu_1:=p(\mathrm{bad})\le4d\ell\le4dC_0/n\le\frac1{16}$.
\item At a triangle-good $\omega$, $U_f(\xi\otimes\ket\omega)=\sum_gA_g\xi\otimes\ket{f(g)\omega}$. Lemma~\ref{lem:restrict-biased} on the good points, a diagonal projection that commutes
with $p$, gives a bath $p_1$ with leakage $0$, distance at most $(\frac1{16}+\frac1{16})/\frac{15}{16}=\frac2{15}$, $S_1\le S/(1-\nu_1)$ and $a_1\le(a+h_2(\nu_1))/(1-\nu_1)$.
\item By Lemma~\ref{lem:perm-leak}, $p_1\{\omega:f(g)\omega=f(h)\omega\}\le2\cdot\frac2{15}=\frac4{15}$ for $g\ne h$.
\item Lemma~\ref{lem:tensor} applied to the round of step 2: on tuples of good points $U_{f^{\times t}}$ has Kraus form with the product model, the bath is $p_1^{\otimes t}$, $S=tS_1$,
$a\le ta_1$, and two labels collide at a tuple if and only if they collide in every coordinate, so the collision mass is at most $(r^2/2)(4/15)^t\le1/n^2$ for
$t=\lceil\log_{15/4}(r^2n^2/2)\rceil$.
\item Lemma~\ref{lem:restrict-biased} again removes the colliding tuples, a diagonal projection with $\nu_2\le1/n^2$. The result is a free exact permutation round, whose
weighted Gram is the identity, so its distance is $0$ (Lemma~\ref{lem:zero-leak}); $S'\le tS_1/(1-\nu_2)$ and $a'\le(ta_1+h_2(\nu_2))/(1-\nu_2)$. Then
$n\,a'\le(1+O(\nu_1))\,t\,n\,a+t\,n\,h_2(4dC_0/n)(1+O(n^{-2}))+O(1)$, and $n\,h_2(4dC_0/n)\le4dC_0\log(en/(4dC_0))=O(dC_0\log n)$.
\end{enumerate}
Theorem~\ref{thm:pointblock} and Lemma~\ref{lem:classperm} need freeness and exactness only at the points that the windows contain, so for the output the windows of the balance
condition must lie inside the good, collision-free tuples, an exit into a bad or colliding tuple counting as a kill.
\end{proof}

\begin{proof}[Proof of Lemma~\ref{lem:conditional}]
Let $\Phi_d$ be maximally entangled on the input and a reference $R_{\rm in}$, and $\Omega:=(U\otimes I)(\ket{\Phi_d}\otimes\ket{\psi_\sigma})$, a pure state of the output, $\C^M$,
$R_{\rm in}$ and $R$. Its $\C^M$-marginal is $T(\sigma)$ and its $R$-marginal has the spectrum of $\sigma$. As $\Omega$ is pure, $S(CR)=S(\C^M)=S(T(\sigma))$, so
$H(C|R)=S(CR)-S(R)=a(\sigma)$. If the round is exact on $K$, then $W=\sum_gA_g\otimes B_g$ there and $\Tr(A_h^*A_g)=d\mu_g\delta_{gh}$, so
$T(\sigma)=\sum_g\mu_gB_g\sigma B_g^*=\widetilde X\widetilde X^*$ with $\widetilde X:=\sum_g\sqrt{\mu_g}B_g\sqrt\sigma\bra g:\C^r\otimes K\to\C^M$. Then
$\widetilde X^*\widetilde X=\sum_{g,h}\sqrt{\mu_g\mu_h}\ket g\bra h\otimes\sqrt\sigma B_g^*B_h\sqrt\sigma$ has the nonzero spectrum of $T(\sigma)$ and $K$-marginal
$\sum_g\mu_g\sqrt\sigma B_g^*B_g\sqrt\sigma=\sigma$, so its conditional entropy given $K$ is $a(\sigma)$.
\end{proof}

We also use the leaky form of this identity. For a state $\sigma$ on $K$ with leakage $\ell_\sigma=\Tr(\sigma\Lambda)/d$ put
$\omega^0:=\sum_{g,h}\sqrt{\mu_g\mu_h}\ket g\bra h\otimes\sqrt\sigma X_g^*X_h\sqrt\sigma$, of trace $1-\ell_\sigma$. Then
\begin{equation}\label{eq:leaky-conditional}
(1-\ell_\sigma)\,H(L|R)_{\omega^0/(1-\ell_\sigma)}\ \le\ a(\sigma)+h_2(\ell_\sigma)+\ell_\sigma\log(d^2-r).
\end{equation}
Indeed, write $\ket{\psi_\sigma}=(\sqrt\sigma\otimes I)\ket\Gamma$ with $\ket\Gamma=\sum_\kappa\ket{\kappa\kappa}$, and $W=G+L$. The part of $\Omega$ coming from $G$ is
$\sum_gd^{-1/2}\vert A_g\rangle\!\rangle\otimes(X_g\sqrt\sigma\otimes I)\ket\Gamma$, and the part coming from $L$ lies in the range of $\Pi^\perp\otimes I$, $\Pi:=\Pi_\Phi\otimes I_R$, by
Hilbert--Schmidt orthogonality. Let $\sigma_{CR}$ be the $CR$-marginal of $\Omega$. Since $\Tr_M[(A\otimes I)\ket\Gamma\bra\Gamma(B\otimes I)^*]=(B^*A)^T$, the block $(g,h)$ of
$\Pi\sigma_{CR}\Pi$ in the orthonormal basis $\vert A_g\rangle\!\rangle/\sqrt{d\mu_g}$ is the entrywise conjugate of the block $(g,h)$ of $\omega^0$; conjugation preserves spectra and
commutes with partial traces, and $\Tr(\Pi\sigma_{CR}\Pi)=1-\ell_\sigma$, so $H(C|R)_{\Pi\sigma_{CR}\Pi/(1-\ell_\sigma)}=H(L|R)_{\omega^0/(1-\ell_\sigma)}=:H_0$. The pinched state
$\tilde\sigma:=\Pi\sigma_{CR}\Pi+\Pi^\perp\sigma_{CR}\Pi^\perp$ has the $R$-marginal of $\sigma_{CR}$. With $A:=\Pi\sqrt{\sigma_{CR}}$, $\Pi\sigma_{CR}\Pi=AA^*$ has the nonzero spectrum of
$\sqrt{\sigma_{CR}}\Pi\sqrt{\sigma_{CR}}$, and $\sigma_{CR}=\sqrt{\sigma_{CR}}\Pi\sqrt{\sigma_{CR}}+\sqrt{\sigma_{CR}}\Pi^\perp\sqrt{\sigma_{CR}}$, so concavity gives the inequality of Groenewold
and Lindblad~\cite{Lindblad1972}, $S(\sigma_{CR})\ge(1-\ell_\sigma)S(\Pi\sigma_{CR}\Pi/(1-\ell_\sigma))+\ell_\sigma S(\Pi^\perp\sigma_{CR}\Pi^\perp/\ell_\sigma)$, while $S(\tilde\sigma)$ is $h_2(\ell_\sigma)$
plus the same average; so $H(C|R)_{\tilde\sigma}\le a(\sigma)+h_2(\ell_\sigma)$. Finally $\tilde\sigma$ is block diagonal for $\Pi\oplus\Pi^\perp$ on $C$, so with the flag $F$ of the
block, $H(C|R)_{\tilde\sigma}=H(F|R)+H(C|FR)\ge(1-\ell_\sigma)H_0-\ell_\sigma\log(d^2-r)$, since the $\Pi^\perp$ block lives on a space of dimension $d^2-r$.

\begin{proof}[Proof of Lemma~\ref{lem:domination}]
From $W^*W=I$, $H:=G^*G-I=-(G^*L+L^*G+L^*L)$. From $(\sqrt tG\pm L/\sqrt t)^*(\sqrt tG\pm L/\sqrt t)\ge0$, $\mp(G^*L+L^*G)\le tG^*G+L^*L/t$. Hence $H\le tG^*G+L^*L/t$ and
$-H\le tG^*G+(1+1/t)L^*L$. The first gives $(1-t)G^*G\le I+L^*L/t$, so $tG^*G\le(tI+L^*L)/(1-t)$, and inserting it,
$\pm H\le[tI+(1+1/t-t)L^*L]/(1-t)\le\bar B:=(tI+(1+1/t)L^*L)/(1-t)$. For $\ket p\in\C^d$ write $Y_{pq}=(\bra p\otimes I)Y(\ket q\otimes I)$. Then
$\bar B_{pp}=(tI+(1+1/t)(L^*L)_{pp})/(1-t)\le\beta_t$, because $(L^*L)_{pp}\ge0$ and $\sum_p(L^*L)_{pp}=\Lambda$. By (F1), $G(\ket{z_g}\otimes\kappa)=\ket{z_g}\otimes X_g\kappa$, so
$H_{z_gz_g}=X_g^*X_g-I$, and $-\beta_t\le H_{z_gz_g}\le\beta_t$. Compress $\bar B+H\ge0$ and $\bar B-H\ge0$ to two slots $p,q$, conjugate the second by $\diag(I,-I)$ and add:
$\left(\begin{smallmatrix}\bar B_{pp}&H_{pq}\\H_{qp}&\bar B_{qq}\end{smallmatrix}\right)\ge0$. For the two columns $p,q$ of a triangle block, $(G^*G)_{pq}=\frac12R_t$ and $I_{pq}=0$; raising
the diagonal blocks to $\beta_t$ keeps positivity, and multiplying by $2$ gives the claim.
\end{proof}

Conjugating the last matrix by $\diag(I,e^{i\phi})$ or by $\diag(I,-I)$, or swapping its blocks, and scaling, shows
$\left(\begin{smallmatrix}2|c|\beta_t&cD\\\bar cD^*&2|c|\beta_t\end{smallmatrix}\right)\ge0$ for every triangle defect or adjoint $D$ and $c\in\C$. Hence, if $E$ is Hermitian on
$\C^r\otimes K$ with zero diagonal blocks and $E_{pq}=\sum_ec_{pq,e}D_e$ with every $D_e$ a triangle defect or its adjoint up to a phase, then
$\pm E+I_r\otimes2\bar c\beta_t\ge0$ with $\bar c:=\max_p\sum_{q\ne p}\sum_e|c_{pq,e}|$: embed one such $2\times2$ block matrix at the positions $(p,q)$ and $(q,p)$ for every $p<q$ and every
term, sum, and add $\bigoplus_p(2\bar c-\sum_{q\ne p}\sum_e2|c_{pq,e}|)\beta_t\ge0$. For the correction $E$ that projects the off-diagonal blocks of $Y^0=[X_g^*X_h]$ onto the
identities (the average over each class of linked label pairs, as in the proof of Theorem~\ref{thm:F}), each block difference telescopes along at most $r(r-1)-1$ edges,
so $\bar c\le(r-1)(r(r-1)-1)\le(r-1)(r^2-1)$.

\begin{proof}[Proof of Theorem~\ref{thm:BF}]
\emph{Facts.} For every state $\sigma$ on $K$, $\ell_\sigma=\Tr(\sigma\Lambda)/d$, as in (f1) of the proof of Theorem~\ref{thm:F}. Applying $\Tr_d$ to $W^*W=I$, the cross terms vanish
and $\Tr_d(G^*G)=\sum_gd\mu_gX_g^*X_g$, so $\sum_g\mu_gX_g^*X_g=I-\Lambda/d$ and $\sum_g\mu_g\Tr(\sigma X_g^*X_g)=1-\ell_\sigma$. For $g\ne h$ and any state $\sigma$ on $K$ with distance
$u_\sigma$,
\begin{equation}\label{eq:anchor}
|\Tr(\sigma X_g^*X_h)|\le2u_\sigma+2\sqrt{\Tr\sigma\Lambda}+\Tr\sigma\Lambda :
\end{equation}
the anchor block of $W$ is $W_{z_gz_g}=X_g+\lambda_g$ with $\lambda_g=\sum_\alpha(E_\alpha)_{z_gz_g}L_\alpha$ and $\lambda_g^*\lambda_g\le\Lambda$ by Cauchy--Schwarz;
$|\Tr(\sigma W_{z_hz_h}^*W_{z_gz_g})|\le2u_\sigma$ by (F1) and (F2); $\nrm{W_{z_gz_g}\sqrt\sigma}_F\le1$ since $W_{zz}^*W_{zz}\le(W^*W)_{zz}=I$; $\nrm{\lambda_g\sqrt\sigma}_F^2\le\Tr\sigma\Lambda$;
expand $X=W_{zz}-\lambda$ and use Cauchy--Schwarz.

\begin{enumerate}[label=\arabic*.,itemsep=1pt]
\item \emph{Commuting restriction.} Write $\rho=\sum_i\rho_i\ket{e_i}\bra{e_i}$ ($\rho_i>0$) and $\lambda_i:=\bra{e_i}\Lambda\ket{e_i}$, so $\sum_i\rho_i\lambda_i=d\ell=s^2$. Put $\theta:=s/c_r$,
$J:=\{i:\lambda_i\le\theta\}$, $P_0:=\sum_{i\in J}\ket{e_i}\bra{e_i}$ and $K_0:=\operatorname{ran}P_0$. Markov's inequality gives $\nu:=\sum_{i\notin J}\rho_i\le s^2/\theta=c_rs\le\frac1{46}$. $P_0$ commutes
with $\rho$, so Lemma~\ref{lem:restrict-biased} gives $\rho_0:=P_0\rho P_0/(1-\nu)$ with $S_0\le S/(1-\nu)$, $a_0\le(a+h_2(\nu))/(1-\nu)$, $u_0\le(u+\nu)/(1-\nu)$ and $\ell_0\le\ell/(1-\nu)$.
On $K_0$, $W_0=G_0+L_0$ with the $X_g$ restricted and $\Lambda_0=P_0\Lambda P_0$. In the basis $\{e_i\}_{i\in J}$, $\rho_0\log(1/\rho_0)$ is diagonal with nonnegative entries, so
$J_0:=\Tr(\Lambda_0\rho_0\log(1/\rho_0))=\sum_{i\in J}\lambda_i\rho_{0,i}\log(1/\rho_{0,i})\le\theta S_0$. From now on everything is on $K_0$, where $\rho_0$ is invertible; $\xi:=\sqrt{\rho_0}$.
\item \emph{Domination.} $Y^0:=[X_g^*X_h]_{g,h}\ge0$ on $\C^r\otimes K_0$, and $\beta:=(sI+(1+1/s)\Lambda_0)/(1-s)$; Lemma~\ref{lem:domination} applies to $W_0$ with $t=s$.
\item \emph{Projection.} Let $E$ be the class-average correction, so that the off-diagonal blocks of $Y^0+E$ satisfy every identity, and $Z_E:=2(r-1)(r^2-1)\beta\ge2\bar c\beta$. Then
$\pm E+I_r\otimes Z_E\ge0$.
\item \emph{Fill.} $Z_D:=\beta$, $Z:=Z_D+Z_E=c_r\beta$, $F:=\bigoplus_g(I+Z_D-X_g^*X_g)\ge0$ by Lemma~\ref{lem:domination}, and $Y_1:=Y^0+F+(E+I_r\otimes Z_E)\ge0$. Then $(Y_1)_{gg}=I+Z$
for every $g$ and $(Y_1)_{gh}=Y^0_{gh}+E_{gh}$ for $g\ne h$, and every identity holds: the off-diagonal ones by step 3, the diagonal ones because all diagonal blocks are equal.
\item \emph{Absorber in the bath.} $z:=\Tr(\rho_0Z)$ and $\rho':=(\rho_0+\xi Z\xi)/(1+z)$, a state on $K_0$ with $\rho'\ge\rho_0/(1+z)$. Put
$Y:=(1+z)^{-1}(I\otimes\rho'^{-1/2}\xi)Y_1(I\otimes\xi\rho'^{-1/2})$. Then $Y\ge0$, $Y_{gg}=\rho'^{-1/2}(\rho_0+\xi Z\xi)\rho'^{-1/2}/(1+z)=I$, and every identity survives the common
congruence. $X'_g:=Y^{1/2}(\ket g\otimes I):K_0\to\C^r\otimes K_0$ are isometries with $X'^*_gX'_h=Y_{gh}$. With the embedding $\iota:=X'_1$, the bath $X'_1\rho'X'^*_1$ and the family
$X'_gX'^*_1$ on $\operatorname{ran}X'_1$, Lemma~\ref{lem:dilation}(c) extends $\sum_gA_g\otimes X'_gX'^*_1$ to a unitary. The new round has channel
$\sum_{g,h}\Tr(\rho'X'^*_hX'_g)A_g(\cdot)A_h^*$, hence leakage $0$, bath entropy $S':=S(\rho')$, and bath map with the spectrum of $\sum_g\mu_gX'_g\rho'X'^*_g$.
\item \emph{Entropy production.} $(I\otimes\sqrt{\rho'})Y(I\otimes\sqrt{\rho'})=Y_1^\xi/(1+z)$ with $Y_1^\xi:=(I\otimes\xi)Y_1(I\otimes\xi)$. By Lemma~\ref{lem:conditional}, $a'=H(L|R)_{\omega'}$ with
$\omega':=DY_1^\xi D/(1+z)$, $D:=\diag(\sqrt{\mu_g})\otimes I$. Let $\omega^0:=D(I\otimes\xi)Y^0(I\otimes\xi)D$, of trace $1-\ell_0$, $\hat\omega^0:=\omega^0/(1-\ell_0)$, and
$\Delta:=D(I\otimes\xi)(F+E+I\otimes Z_E)(I\otimes\xi)D\ge0$. As $E$ has zero diagonal blocks, $\Tr\Delta=\sum_g\mu_g\Tr(\rho_0(I+Z_D-X_g^*X_g))+\Tr(\rho_0Z_E)=\ell_0+z$. Hence
\[
\omega'=(1-p)\hat\omega^0+p\,\tau,\qquad p:=(\ell_0+z)/(1+z),\quad\tau:=\Delta/(\ell_0+z):
\]
a mixture with a positive correction. The mixture bound ($H(L|R)\le H(LX|R)\le h_2(p)+(1-p)H(L|R)_{\hat\omega^0}+p\,H(L|R)_\tau$ for an adjoined index $X$) and $H(L|R)_\tau\le\log r$ give
$a'\le(1-p)H_0+p\log r+h_2(p)\le H_0+2p\log r+h_2(p)$ with $H_0:=H(L|R)_{\hat\omega^0}\ge-\log r$. By~\eqref{eq:leaky-conditional} for $(U,\rho_0)$,
$(1-\ell_0)H_0\le a_0+h_2(\ell_0)+\ell_0\log(d^2-r)$; with step 1 this is item 3.
\item \emph{Bath entropy.} $\log$ is operator monotone and $\rho'\ge\rho_0/(1+z)$, so $S(\rho')\le-\Tr\rho'\log\rho_0+\log(1+z)=(S_0+\Tr(\xi Z\xi\log(1/\rho_0)))/(1+z)+\log(1+z)$. Since $\xi$
commutes with $\log\rho_0$, $\Tr(\xi Z\xi\log(1/\rho_0))=\Tr(Z\rho_0\log(1/\rho_0))=\frac{c_r}{1-s}(sS_0+(1+1/s)J_0)\le\frac{(c_rs+s+1)S_0}{1-s}$ by step 1, as $\theta=s/c_r$. So
$S'\le S_0(1+(1+(c_r+1)s)/(1-s))+z/\ln2$.
\item \emph{Gram.} For $g\ne h$, $\Tr(\rho'X'^*_gX'_h)=\Tr((Y_1^\xi)_{gh})/(1+z)=[\Tr(\rho_0X_g^*X_h)+\Tr(\xi E_{gh}\xi)]/(1+z)$. Compressing $\pm E+I\otimes Z_E\ge0$ to the blocks $g,h$ and
conjugating by $\xi$ gives $\left(\begin{smallmatrix}\xi Z_E\xi&\pm\xi E_{gh}\xi\\\pm\xi E_{hg}\xi&\xi Z_E\xi\end{smallmatrix}\right)\ge0$; a matrix
$\left(\begin{smallmatrix}A&X\\X^*&A\end{smallmatrix}\right)\ge0$ has $X=A^{1/2}CA^{1/2}$ with $\nrm C\le1$, so $\nrm{\xi E_{gh}\xi}_1\le\Tr(\rho_0Z_E)$. With~\eqref{eq:anchor} for $\rho_0$,
$|\Tr(\rho'X'^*_gX'_h)|\le2u_0+2\sqrt{d\ell_0}+d\ell_0+\Tr(\rho_0Z_E)$.
\item \emph{Constants.} $c_r\ge7$, so $s\le1/322$ and $\nu\le c_rs\le\frac1{46}$; $d\ell_0\le s^2/(1-\nu)$. Then $z=c_r(s+(1+1/s)d\ell_0)/(1-s)\le c_rs(1+(1+s)/(1-\nu))/(1-s)\le2.05c_rs$,
$\Tr(\rho_0Z_E)=(c_r-1)z/c_r\le2.05(c_r-1)s$ and $p\le\ell_0+z\le2.2c_rs$. For the Gram, $2(u+\nu)/(1-\nu)\le2u(1+1.05c_rs)+2.05c_rs$ and $2\sqrt{d\ell_0}+d\ell_0\le2.1s$, so the bound is
at most $2u(1+1.05c_rs)+4.2c_rs$, which at $u\le\frac1{16}$ and $c_rs\le\frac1{46}$ is at most $0.128+0.092<\frac14$. For the bath, $(1+(c_r+1)s)/(1-s)\le1+1.3c_rs$ and
$S_0\le S(1+1.03c_rs)$, so $S'\le(2+1.3c_rs)(1+1.03c_rs)S+2.05c_rs/\ln2\le(2+3.4c_rs)S+3.1c_rs$. At the scale of the law, $s=\sqrt{d\ell}\le\sqrt{dc_1/n}\le1/(46c_r)$ when
$n\ge2116c_r^2dc_1$, and then $c_rs\le\frac1{46}$ gives $S'\le2.074S+0.068\le2.08S+0.07$. In item 3, $1/((1-\nu)(1-\ell_0))\le1+1.1c_rs$, $h_2$ is increasing on $[0,\frac12]$,
$n\cdot2p\log r\le4.4c_r\sqrt{dc_1n}\log r$, $n\,h_2(\ell_0)/(1-\ell_0)\le1.1c_1\log(en/c_1)$ and $n\ell_0\log(d^2-r)/(1-\ell_0)\le2.2c_1\log d$; and $n\,h_2(x)\le nx\log(e/x)=O(\sqrt n\log n)$
for $x=O(n^{-1/2})$. This gives $n\,a'\le(1+O(n^{-1/2}))\,n\,a+O_{d,r,c_1}(\sqrt n\log n)$.\qedhere
\end{enumerate}
\end{proof}

\begin{proof}[Proof of Corollary~\ref{cor:BF}]
The lower inequality, with these constants, is Theorem~\ref{thm:lower-constants}. For the upper one take a round with $u\le\frac1{16}$, $\ell\le\frac2{15n}$ and $S+n\,a\le\Ls(n)+1$. Theorem~\ref{thm:BF} with $c_1=\frac2{15}$
gives a zero-leakage round with weighted Gram off-diagonals at most $\frac14$, $S'\le2.08S+0.07$ and $n\,a'\le(1+O(n^{-1/2}))\,n\,a+O(\sqrt n\log n)$, so
$S'+n\,a'\le2.1(S+n\,a)+O(\sqrt n\log n)$ for large $n$. Apply Theorem~\ref{thm:coherent}, which gives the displayed squared-log upper with flat imports and charged initial purity.
\end{proof}
